\SourceLeaf{387}%
Heft (1810) wurde sie ganz eingestellt. General \textsc{Guilleminot}, der 1822 die Leitung des Depots übernahm, fasste als Erster den Plan, die Zeitschrift regelmäßig fortzusetzen. Sein kommissarischer Stellvertreter, General \textsc{Delachasse de Verigny}, arbeitete den Plan bis zur Ausführungsreife aus. Zugleich beschloss man, die Form etwas zu ändern und auch die sieben älteren Hefte. deren Exemplare vergriffen waren, in dem für die Fortsetzung gewählten Quartformat neu zu drucken. Für diesen Neudruck der früheren Hefte sind die ersten beiden Bände der neuen Ausgabe vorgesehen. Der erste liegt uns zusammen mit drei Bänden der Fortsetzung vor, der zweite soll demnächst nachgeliefert werden.

Das \textit{Dépôt de la guerre} wurde erstmals 1688 unter der Regierung \textsc{Ludwigs} XIV. vom Minister \textsc{Louvois} gegründet. Anfangs war es jedoch nur ein Archiv, in dem der gesamte amtliche Schriftverkehr der Armee hinterlegt wurde. Seine heutige Gestalt gewann es erst nach und nach, als sich sein Aufgabenbereich erweiterte und seine Organisation festigte. Seinen besonderen, großartigen Charakter erhielt es erst, seit es zur zentralen Stelle geworden ist, an der alle Ergebnisse der Arbeit eines selbstständigen Korps, der Ingenieurgeografen, zusammenlaufen. Einen Eindruck vom Umfang der Tätigkeit der Einrichtung gibt die Tatsache, dass ihre jährlichen Kosten 1801 auf 110000 Franken veranschlagt wurden. Darin waren die Gehälter des Personals noch nicht enthalten. Diese beliefen sich 1793 auf 231100 Franken.

An eine Zeitschrift, die aus einer so überaus reichen Quelle schöpfen kann, darf man hohe Ansprüche stellen. Sie werden umso besser erfüllt, je mehr die Herausgeber darauf achten, wichtiges, der Öffentlichkeit bisher unzugängliches Material zu veröffentlichen und alles auszuschließen, was die Wissenschaft nicht weiterbringt.

Wir finden die älteren Artikel viel weiter von diesem Ideal entfernt als die neuere Fortsetzung. Abgesehen von einem Aufsatz über die Hydrographie eines Teils von Frankreich und einer Notiz zur Geschichte des \textit{Dépôt de la guerre} hätten die Artikel des ersten Bandes tatsächlich ebenso gut geschrieben werden können, wenn es das \textit{Dépôt} gar nicht gegeben hätte. Ohne das Interesse zu bestreiten, das mehrere Aufsätze vor dreißig Jahren wecken konnten und zum Teil noch heute wecken, kann
